\section{統計環境Rの導入}
\label{s03_introR}
\subsection{授業内容}
\subsubsection{科目の中でのこのコマの位置づけ}
\begin{tabular}[t]{ccccccc}
    記述統計[0] & 確率[0] & モデル[0] & 推測・推定[0] & 検定[0] & 実践[1]
\end{tabular}
\subsubsection{概要}
携帯電話やタブレットが誰でも身近に利用できるようになってきてはいるが，それらはあくまでも受動的にサービスを利用する媒体であり，執筆や計算などユーザが能動的に電子機器を利用するにあたっては，パソコンの利用は不可欠である。
本講義だけでなく，心理学の研究をする上でも当然パソコンの利用は必須であり，特に統計的データ処理をするにあたってはより深く知る必要がある。また，ExcelやNumbersなどの表計算でも記述統計量やグラフの作図はできるが，それ以上の「複数の変数を扱った多変量解析」や「群間の平均値の差を細かく検証する」といった心理統計的な作業には不十分なのである。
そこで，より専門的なツールとして統計環境Rを導入する。統計環境Rは専修大学人間科学部心理学科が公式に採用している統計環境であり，これを活用する総合開発環境であるRStudioの基本的な使い方，結果の再現性を担保するための工夫などについても言及する。それに先立って，コンピュータの基礎知識についても解説する。
\subsubsection{コマ主題細目}
\begin{description}

    \item[パソコンのファイル構造]あらゆる媒体が情報という一元化された単位で扱われるのは便利なことではあるが，一見してその違いがわからないのは不便である。そこでファイル名には拡張子と呼ばれるファイルの種類を区別する情報を付与する。また，ファイルがたくさん出来上がって管理できないのも困るので，いくつかのファイルをまとめるフォルダ/ディレクトリという単位でまとめて置くことを理解する。またこれらのファイルがどこにあるかをしっかり理解する必要がある。特にクラウド/ローカルの違いについて注意しておく。

    \item[統計環境R] 心理学科では基本的に統計環境Rを共通のツールとして用いる。統計環境Rはフリーソフトウェアであり，誰でも無償で利用できる環境である。また多くのパッケージを用いることで最先端の分析を行うこともでき，基本的にRでできない統計解析はない，と考えてもらって間違いない。まずは統計環境Rの基本的な知識である，対話的な統計処理の方法について，また電卓や表計算ソフトと違って，大量のデータを統計的に処理することについての基本的な理解をする。
    \begin{flushright}
        $\to$ Rについての入門書は多くあるがここでは\citet{kosugi2019}のP.1--7を挙げておく。
    \end{flushright}
    \item[総合開発環境RStudio]統計環境Rは統計言語とも言われ，本体単体では簡素な言語的やりとりができるに過ぎない。ファイルの操作や描画機能，記録，参照，管理など統計言語を扱うに当たって必要な総合的環境としてRStudioを併用することで，その利便性は飛躍的に向上する。まずはRStudioの基本的な使い方と，プロジェクト管理について解説する。
    \begin{flushright}
        $\to$ 同じくRStudioを使った入門書は多くあるがここでは\citet{kosugi2019}のP.7--24を挙げておく。
    \end{flushright}

    \item[再現性のある研究実践へ]今後，パソコンを使って文章のみならず，図表も作成することになる。統計的分析結果の出力も，プログラムによって行うことになる。コンピュータにおいてプログラムが実行されるとはどういうことか(適切なファイルへのアクセス)を理解し，またプログラムを書いて実行するときの注意点について理解する。特にエラーとその報告の仕方については，恐怖心に駆られることなく対応するべきである。そして研究の実践とその成果が「動作」ではなく「指示書」として記録に残ることは，研究の再現性を高めるためにも重要な営みであることを理解する。Gitによる文書のバージョン管理や，OSFなど情報の共有など，最新の研究実践方法についても言及する。


\end{description}
\subsubsection{キーワード}
\begin{itemize}
    \item ファイル，フォルダ，拡張子
    \item RとRStudio
    \item プロジェクト
    \item 再現性のある研究実践
\end{itemize}
\subsection{授業情報}
\paragraph{コマの展開方法}
講義
\paragraph{標準シラバスにおける位置づけ}
科目番号5；心理学統計法；2.統計に関する基礎的な知識；C エクセル，R，SPSS入門(１) 統計分析のための言語の手ほどき　プログラムの基本的操作
\subsubsection{予習・復習課題}
\paragraph{予習}大学生活ではPCを使うことがほぼ確実に必要であり，今後の社会においてもその能力は間違いなく生きてくるので，なるべく早い段階でPCを所持し，ブラインドタッチなど基本的な技術を身につけておく。
\paragraph{復習}自分のPCのユーザ名，ファイル名，拡張子の付け方などを確認し，半角英数になっていないか確認しよう。今後ファイルを作ることがあれば，半角英数文字で名前をつけるよう癖づけておいたほうが良い。これらは未然にエラーを防ぐことにもなる。
